\chapter{Categoría de Baire y principio de contracción}\label{cap:nuevo04}\label{chap:10}\label{chap:11}
La completitud permite convertir procesos infinitos en resultados de existencia. Estudiaremos dos manifestaciones: el teorema de Baire, que hace compatibles familias numerables de condiciones abiertas y densas, y el principio de contracción, que construye una solución de $x=T(x)$ mediante iteraciones y controla su error. Comenzaremos con bolas cerradas anidadas cuyos radios tienden a cero; después desarrollaremos las consecuencias topológicas de Baire y las aplicaciones del punto fijo a sistemas, ecuaciones integrales y conjuntos autosimilares. En ambas partes será esencial justificar que los límites permanecen en los conjuntos donde se imponen las condiciones.

\section{De bolas cada vez menores a un punto}
\label{lectura:10}

\begin{teorema}[Bolas cerradas anidadas]\label{ch10:bolas}
Sea $X$ completo y sea $\overline B(x_n,r_n)$ una sucesión de bolas cerradas no vacías con
\[
 \overline B(x_{n+1},r_{n+1})\subset\overline B(x_n,r_n),
 \qquad r_n\longrightarrow0.
\]
Su intersección contiene exactamente un punto.
\end{teorema}
\begin{proof}
Para $m,k\geq n$, ambos centros pertenecen a $\overline B(x_n,r_n)$, y por tanto $d(x_m,x_k)\leq2r_n$. Esto prueba que $(x_n)$ es de Cauchy. Sea $x$ su límite. Fijado $n$, toda la cola $x_m$ con $m\geq n$ está en la bola cerrada número $n$; su límite también pertenece a ella. Como $n$ era arbitrario, $x$ pertenece a la intersección. Si $y$ perteneciera también, tendríamos $d(x,y)\leq2r_n$ para todo $n$, y dejar $n\to\infty$ da $x=y$.
\end{proof}
La bola cerrada $\overline B(x,r)=\{y:d(x,y)\leq r\}$ no debe confundirse automáticamente con la clausura de $B(x,r)$: en un espacio métrico general pueden ser diferentes. La construcción siguiente usará directamente bolas cerradas.

\section{Conjuntos pequeños en sentido topológico}
\begin{definicion}\index{F sigma@$F_\sigma$}\index{G delta@$G_\delta$}
Un conjunto es de tipo $F_\sigma$ si es unión numerable de cerrados, y de tipo $G_\delta$ si es intersección numerable de abiertos. Por las leyes de De Morgan, el complemento de un $F_\sigma$ es un $G_\delta$, y recíprocamente. Estas denominaciones se usarán para describir conjuntos de continuidad y de discontinuidad.
\end{definicion}


\begin{definicion}
Un subconjunto $A\subset X$ es \emph{denso en ninguna parte} si $\interior{\cl A}=\varnothing$. Es \emph{de primera categoría}, o \emph{magro}, si es unión numerable de conjuntos densos en ninguna parte. Es \emph{de segunda categoría} si no es de primera categoría. Un conjunto cuyo complemento es magro se llama \emph{residual}.
\end{definicion}
\index{Baire, teorema de}\index{categoria de Baire@categoría de Baire}\index{denso en ninguna parte}\index{conjunto!magro}\index{conjunto!residual}
El uso de la clausura es esencial. Los racionales tienen interior vacío en $\R$, pero su clausura es $\R$ y por ello no son densos en ninguna parte. Sí son magros: son unión numerable de conjuntos unipuntuales, cada uno cerrado y de interior vacío.

\begin{ejemplo}
En un espacio discreto, todo subconjunto es abierto y cerrado; por tanto, el único conjunto denso en ninguna parte es el vacío. Todo subconjunto no vacío es de segunda categoría. Así, $\Z$ es magro dentro de $\R$, porque es numerable y cada punto tiene interior vacío, pero es de segunda categoría dentro de sí mismo con la topología heredada, que es discreta. La categoría siempre se refiere a un espacio ambiente concreto.
\end{ejemplo}

\begin{teorema}[Baire, forma densa]\label{ch10:baire}
En un espacio métrico completo, la intersección de una familia numerable de abiertos densos es densa. Equivalentemente, el complemento de todo conjunto magro es denso.
\end{teorema}
\begin{proof}
Sean $G_1,G_2,\ldots$ abiertos densos, y sea $U$ un abierto no vacío. Vamos a construir un punto de $U\cap\bigcap_nG_n$. La densidad de $G_1$ hace no vacío al abierto $U\cap G_1$. Elegimos $x_1$ en él y $r_1>0$ suficientemente pequeño para que
\[
 \overline B(x_1,r_1)\subset U\cap G_1,\qquad r_1<1/2.
\]
La elección es posible porque, si $B(x_1,s)\subset U\cap G_1$, cualquier $0<r_1<s$ sirve. Suponiendo construida la bola número $n$, el conjunto $B(x_n,r_n)\cap G_{n+1}$ es abierto no vacío. Elegimos en él un punto $x_{n+1}$ y un radio $r_{n+1}>0$ con
\[
 \overline B(x_{n+1},r_{n+1})\subset B(x_n,r_n)\cap G_{n+1},
 \qquad r_{n+1}<2^{-(n+1)}.
\]
Las bolas cerradas son anidadas y sus radios tienden a cero. Por el teorema~\ref{ch10:bolas}, tienen un punto común $x$. La primera bola asegura $x\in U$, y la bola número $n$ asegura $x\in G_n$. Esto prueba densidad.

Si $A=\bigcup_nA_n$ con $\interior{\cl{A_n}}=\varnothing$, entonces $G_n=X\setminus\cl{A_n}$ es abierto y denso. La intersección de los $G_n$ está contenida en $X\setminus A$, que también será denso. A la inversa, para abiertos densos $G_n$, cada complemento $X\setminus G_n$ es cerrado de interior vacío, y el complemento de su unión es $\bigcap_nG_n$. Las dos formulaciones son, pues, equivalentes.
\end{proof}

\begin{corolario}\label{ch10:segunda}
Todo espacio métrico completo no vacío es de segunda categoría en sí mismo. Todo espacio métrico completo no vacío y sin puntos aislados es incontable.
\end{corolario}
\begin{proof}
Si $X$ fuera magro en sí mismo, el teorema daría que su complemento vacío es denso, imposible cuando $X\ne\varnothing$. Si además no tiene puntos aislados, cada singleton $\{x\}$ es cerrado de interior vacío. Si $X$ fuera numerable, sería unión numerable de esos conjuntos y resultaría magro, contradiciendo la primera afirmación.
\end{proof}

\begin{observacion}[Qué elecciones utiliza la demostración]
La construcción escoge una nueva bola a partir de la anterior. Trabajamos con el marco usual de teoría de conjuntos con elección, que permite estas elecciones sucesivas. En un espacio separable con una base numerable ya enumerada, se puede tomar en cada paso la primera bola básica que satisfaga las inclusiones y la cota de radio; así la regla queda determinada. La discusión de la equivalencia lógica entre versiones del teorema y principios de elección dependiente pertenece a los fundamentos de la teoría de conjuntos y no interviene en ninguna demostración del volumen. Lo que necesitaremos en análisis es el enunciado métrico que acabamos de demostrar.
\end{observacion}

\section{Ejercicios sobre categoría, cardinalidad y discontinuidades}
\begin{ejercicio}[Lectura 10, ejercicio 1]\label{ej:lectura-10-1}
Determine si es necesario que las bolas del teorema de bolas anidadas sean cerradas.
\end{ejercicio}


\begin{ejercicio}[Lectura 10, ejercicio 2]\label{ej:lectura-10-2}
Enuncie y pruebe el recíproco del teorema de bolas cerradas anidadas.
\end{ejercicio}


\begin{ejercicio}[Lectura 10, ejercicio 3]\label{ej:lectura-10-3}
Pruebe que $X$ es completo si y sólo si toda sucesión de cerrados \emph{no vacíos} $A_{n+1}\subset A_n$ con $\diam(A_n)\to0$ tiene intersección no vacía.
\end{ejercicio}


\begin{ejercicio}[Lectura 10, ejercicio 4]\label{ej:lectura-10-4}
Pruebe que $A$ es denso en ninguna parte si y sólo si $X\setminus\cl A$ es denso en $X$.
\end{ejercicio}


\begin{ejercicio}[Lectura 10, ejercicio 5]\label{ej:lectura-10-5}
Demuestre que todo subconjunto de un conjunto magro es magro.
\end{ejercicio}


\begin{ejercicio}[Lectura 10, ejercicio 6]\label{ej:lectura-10-6}
Sea $X$ completo no vacío y sea $X=\bigcup_{n\geq1}F_n$, donde los $F_n$ son cerrados. Pruebe que alguno tiene interior no vacío.
\end{ejercicio}


\begin{ejercicio}[Lectura 10, ejercicio 7]\label{ej:lectura-10-7}
Demuestre que un espacio de Banach de dimensión infinita posee una base algebraica y que toda base algebraica suya es incontable.
\end{ejercicio}


\begin{ejercicio}[Lectura 10, ejercicio 8]\label{ej:lectura-10-8}
Sea $(q_n)$ una enumeración de $\Q$, y para $\eps>0$ defina
\[
 A_\eps=\bigcup_{n\geq1}(q_n-\eps2^{-n-1},q_n+\eps2^{-n-1}).
\]
Determine si $E=\bigcap_{\eps>0}A_\eps$ contiene puntos irracionales además de los racionales.
\end{ejercicio}


\begin{ejercicio}[Lectura 10, ejercicio 9]\label{ej:lectura-10-9}
\begin{enumerate}[label=(\roman*)]
\item Examine la afirmación original: «No existe una función $f:\R\to\R$ discontinua en todo número irracional». Formule y demuestre la versión correcta referente al conjunto exacto de discontinuidades.
\item Para cada conjunto a lo sumo numerable $A\subset\R$, construya una función cuyas discontinuidades sean exactamente los puntos de $A$.
\end{enumerate}
\end{ejercicio}


\begin{ejercicio}[Lectura 10, ejercicio 10]\label{ej:lectura-10-10}
Construya funciones $\R\to\R$ discontinuas en todos los puntos de: \begin{enumerate}[label=(\roman*)]\item $\Q$;\item el conjunto de Cantor.\end{enumerate}
\end{ejercicio}


\section{Resolver una ecuación mediante iteraciones}
\label{lectura:11}
Baire permite satisfacer simultáneamente condiciones abiertas y densas. Esta herramienta explica la incontabilidad de ciertos conjuntos, impide bases algebraicas numerables en espacios de Banach infinitodimensionales y controla los conjuntos de discontinuidad; más adelante sostendrá los grandes principios sobre operadores lineales. Pasamos ahora a una segunda aplicación de la completitud, centrada en construir una solución a partir de una sucesión de aproximaciones.

Una ecuación puede escribirse en la forma $x=T(x)$. Si aplicar $T$ reduce las distancias por un factor fijo menor que uno, repetir la operación produce aproximaciones cada vez más estables. La completitud garantiza que esas aproximaciones desembocan en un punto del espacio. Además de existencia y unicidad, el argumento proporciona cotas explícitas del error.

\begin{definicion}
Una \emph{contracción} de un espacio métrico $X$ es una aplicación $T:X\to X$ para la que existe $0\leq q<1$ con
\[
 d(Tx,Ty)\leq qd(x,y)\qquad(x,y\in X).
\]
Un \emph{punto fijo} es un punto $x_*$ que satisface $T(x_*)=x_*$.
\end{definicion}
\index{contraccion@contracción}\index{punto fijo!teorema de Banach}
El caso $q=0$ incluye las aplicaciones constantes. La desigualdad estricta $d(Tx,Ty)<d(x,y)$ para $x\ne y$ no proporciona necesariamente una constante uniforme $q<1$; el teorema siguiente exige dicha constante.

\begin{teorema}[Punto fijo de Banach]\label{ch11:banach}
Si $X$ es completo y no vacío y $T:X\to X$ es una contracción con constante $q<1$, entonces existe un único punto fijo $x_*$. Para cualquier $x_0\in X$, las iteraciones $x_{n+1}=T(x_n)$ convergen a él. Si $0<q<1$, se tienen las cotas
\[
 d(x_n,x_*)\leq\frac{q^n}{1-q}d(x_1,x_0),
 \qquad d(x_n,x_*)\leq\frac{q}{1-q}d(x_n,x_{n-1})\quad(n\geq1).
\]
\end{teorema}
\begin{proof}
Si $q=0$, $T$ es constante, digamos $T(x)=c$, y $T(c)=c$; además ningún otro punto puede ser fijo. Supongamos $0<q<1$. Por inducción,
\[
 d(x_{n+1},x_n)\leq q^nd(x_1,x_0).
\]
Para $m>n$, sumamos distancias entre términos consecutivos:
\[
 d(x_m,x_n)\leq\sum_{k=n}^{m-1}d(x_{k+1},x_k)
 \leq d(x_1,x_0)\sum_{k=n}^{m-1}q^k
 \leq\frac{q^n}{1-q}d(x_1,x_0).
\]
El extremo derecho tiende a cero y no depende de $m$, así que $(x_n)$ es de Cauchy. Sea $x_*$ su límite en $X$. Una contracción es Lipschitz y continua; luego $T(x_n)\to T(x_*)$. Pero $T(x_n)=x_{n+1}\to x_*$, y la unicidad del límite da $T(x_*)=x_*$. Si $y_*$ fuese otro punto fijo,
\[
 d(x_*,y_*)=d(Tx_*,Ty_*)\leq qd(x_*,y_*),
\]
de donde $(1-q)d(x_*,y_*)\leq0$. Como $1-q>0$, los puntos coinciden.

Pasando $m\to\infty$ en la estimación de las colas obtenemos la primera cota. Para la segunda repetimos la suma desde $x_n$ y usamos $d(x_{n+j+1},x_{n+j})\leq q^{j+1}d(x_n,x_{n-1})$; la serie geométrica suma $q/(1-q)$. Esta última cota permite estimar el error a partir de dos iteraciones ya calculadas.
\end{proof}

\begin{ejemplo}[Existencia por continuidad y existencia por contracción]
Toda función continua $f:[0,1]\to[0,1]$ tiene al menos un punto fijo. En efecto, $h(x)=x-f(x)$ satisface $h(0)\leq0$ y $h(1)\geq0$; el teorema del valor intermedio proporciona un cero. La unicidad no se deduce: la identidad tiene todos los puntos fijos. Por otro lado, si $f:\R\to\R$ es derivable y $|f'(x)|\leq q<1$ para todo $x$, el valor medio da $|f(x)-f(y)|\leq q|x-y|$. Banach produce entonces un único punto fijo y convergencia desde cualquier valor inicial.
\end{ejemplo}

\section{Sistemas no lineales y elección de una norma}
\begin{proposicion}\label{ch11:jacobiano}
Sea $R\subset\R^n$ un rectángulo cerrado no vacío y $T:R\to R$ una aplicación de clase $C^1$ en una vecindad de $R$. Si
\[
 \sup_{x\in R}\max_i\sum_{j=1}^n
 \left|\frac{\partial T_i}{\partial x_j}(x)\right|\leq q<1,
\]
entonces $T$ es una contracción para la norma máxima y tiene un único punto fijo en $R$.
\end{proposicion}
\begin{proof}
El segmento entre $x,y\in R$ permanece en $R$ por convexidad. Aplicando el teorema fundamental del cálculo a $t\mapsto T_i(y+t(x-y))$,
\[
 T_i(x)-T_i(y)=\int_0^1\sum_j
 \frac{\partial T_i}{\partial x_j}(y+t(x-y))(x_j-y_j)\,dt.
\]
El valor absoluto está acotado por $q\norm{x-y}_\infty$. Tomar el máximo en $i$ demuestra la contracción. El rectángulo es cerrado en el espacio completo $\R^n$, por lo que es completo; ahora se aplica el teorema~\ref{ch11:banach}. Además de las derivadas, debe comprobarse que $T(R)\subset R$. Sin esta condición las iteraciones podrían salir del dominio.
\end{proof}
En particular, basta que cada derivada parcial tenga valor absoluto no mayor que $q/n$. La hipótesis de rectángulo cerrado corrige la posible pérdida de completitud en un rectángulo abierto.

\begin{ejemplo}[Un sistema de tres ecuaciones]\label{ch11:sistema3}
Consideremos
\begin{align*}
 3x-\cos(yz)-\tfrac12&=0,\\
 x^2-81(y+0.1)^2+\sin z+1.06&=0,\\
 e^{-xy}+20z+\tfrac{10\pi-3}{3}&=0.
\end{align*}
La rama $y+0.1\geq0$ conduce a
\[
 T(x,y,z)=\left(\frac{\cos(yz)}3+\frac16,
 \frac{\sqrt{x^2+\sin z+1.06}}9-0.1,
 -\frac{e^{-xy}}{20}-\frac{10\pi-3}{60}\right).
\]
Verificamos todas las condiciones en $R=[-1,1]^3$. El radicando $Q=x^2+\sin z+1.06$ satisface
\[
 Q\geq1.06-\sin1>0.21,\qquad Q<2.91<4.
\]
La raíz está bien definida y es diferenciable en una vecindad del cubo. La primera coordenada de $T$ está entre $-1/6$ y $1/2$; la segunda está entre $-0.1$ y $2/9-0.1$, y la tercera satisface
\[
 |T_3|\leq e/20+(10\pi-3)/60<1.
\]
Así $T(R)\subset R$.

Las derivadas no nulas son
\begin{align*}
 \partial_yT_1&=-z\sin(yz)/3,&\partial_zT_1&=-y\sin(yz)/3,\\
 \partial_xT_2&=x/(9\sqrt Q),&\partial_zT_2&=\cos z/(18\sqrt Q),\\
 \partial_xT_3&=ye^{-xy}/20,&\partial_yT_3&=xe^{-xy}/20.
\end{align*}
Las de $T_1$ están acotadas por $\sin1/3<0.3$; las de $T_2$ por $1/(9\sqrt{0.21})<0.3$ y $1/(18\sqrt{0.21})<0.3$; las de $T_3$ por $e/20<0.3$. Cada fila tiene suma menor que $0.9$, de modo que $q=0.9$ es una constante válida. Banach garantiza un único punto fijo de esta aplicación en el cubo.

En este caso el punto fijo puede darse exactamente:
\[
 (x_*,y_*,z_*)=(1/2,0,-\pi/6).
\]
La primera coordenada de $T$ es $1/3+1/6=1/2$; para la segunda el radicando vale $1/4-1/2+1.06=0.81$, cuya raíz es $0.9$; la tercera es $-1/20-(10\pi-3)/60=-\pi/6$. Se comprueban así las tres igualdades. La unicidad afirmada corresponde a la aplicación y a la rama elegida; despejar la raíz positiva no permite descartar sin examen las soluciones de la otra rama del sistema original.
\end{ejemplo}

\section{La distancia de Hausdorff y un fractal como punto fijo}
Sea $\mathcal F(X)$ la familia de subconjuntos \emph{no vacíos}, cerrados y acotados de $X$. Para $A,B\in\mathcal F(X)$ definimos
\[
 H(A,B)=\max\left\{\sup_{a\in A}\dist(a,B),
                         \sup_{b\in B}\dist(b,A)\right\}.
\]
\index{Hausdorff, distancia de}\index{fractal}\index{Cantor, conjunto de}
La no vacuidad evita distancias indefinidas y supremos sobre familias vacías. La acotación hace finitos los dos supremos: fijando $b_0\in B$, por ejemplo, $\dist(a,B)\leq d(a,b_0)$ queda acotada sobre $A$.

\begin{proposicion}\label{ch11:hausdorff-metrica}
La función $H$ es una métrica en $\mathcal F(X)$. Además,
\[
 |\dist(x,A)-\dist(x,B)|\leq H(A,B)\qquad(x\in X).
\]
\end{proposicion}
\begin{proof}
La simetría y la no negatividad están en la definición. Si $H(A,B)=0$, cada $a\in A$ tiene distancia cero a $B$ y, como $B$ es cerrado, pertenece a $B$. Al intercambiar los conjuntos obtenemos $A=B$. La vuelta es inmediata.

Para la desigualdad triangular observamos que, dados $a\in A$ y $\eta>0$, podemos elegir $b\in B$ con $d(a,b)<\dist(a,B)+\eta$ y después $c\in C$ con $d(b,c)<\dist(b,C)+\eta$. Entonces
\[
 \dist(a,C)\leq\dist(a,B)+H(B,C)+2\eta.
\]
Dejando $\eta\downarrow0$ y tomando el supremo en $a$ queda una cota por $H(A,B)+H(B,C)$. La otra distancia dirigida, de $C$ hacia $A$, se acota escogiendo primero un punto de $B$ próximo a un punto de $C$ y después uno de $A$ próximo a ese punto de $B$; da la misma suma. Tomar el máximo prueba $H(A,C)\leq H(A,B)+H(B,C)$.

Para la última afirmación, elegimos $a\in A$ con $d(x,a)<\dist(x,A)+\eta$ y un punto de $B$ a distancia menor que $H(A,B)+\eta$ de $a$. Resulta $\dist(x,B)\leq\dist(x,A)+H(A,B)+2\eta$. Al dejar $\eta\downarrow0$ e intercambiar $A,B$ obtenemos la estimación.
\end{proof}

El ejercicio~\ref{ej:lectura-11-5} demostrará que $\mathcal F(X)$ es completo exactamente cuando $X$ lo es. Para $X=\R$, definamos
\[
 \mathcal T(A)=\{a/3:a\in A\}\cup\{(a+2)/3:a\in A\}.
\]
Las transformaciones afines son homeomorfismos de $\R$, por lo que sus imágenes de cerrados son cerradas; también preservan no vacuidad y acotación. Así $\mathcal T$ actúa en $\mathcal F(\R)$. Para un punto $a/3$ de la primera pieza y un punto $b\in B$ próximo a $a$,
\[
 \dist(a/3,\mathcal T(B))\leq|a-b|/3.
\]
Tomando el ínfimo sobre $b$ se obtiene una cota $\dist(a,B)/3$. Para $(a+2)/3$ se usa la segunda pieza de $\mathcal T(B)$ y resulta la misma cota. Tomando supremos, y después intercambiando $A,B$, concluimos
\[
 H(\mathcal T(A),\mathcal T(B))\leq\tfrac13H(A,B).
\]
Por Banach existe un único conjunto fijo. Es el conjunto de Cantor: en su representación ternaria con dígitos $0$ y $2$, la primera pieza antepone el dígito $0$ y la segunda antepone el dígito $2$. Toda sucesión de esos dígitos comienza con uno de ellos, por lo que $C=C/3\cup(C+2)/3$. Las iteraciones desde cualquier cerrado acotado no vacío convergen a $C$ en distancia de Hausdorff.

\section{Ejercicios sobre contracciones, iteraciones y distancia de Hausdorff}
\begin{ejercicio}[Lectura 11, ejercicio 1]\label{ej:lectura-11-1}
Si $X$ es completo no vacío y $f:X\to X$ satisface que $f\circ f$ es una contracción, determine si $f$ tiene un punto fijo.
\end{ejercicio}


\begin{ejercicio}[Lectura 11, ejercicio 2]\label{ej:lectura-11-2}
Usando el ejercicio~\ref{ej:lectura-11-1}, pruebe que $x=\cos x$ tiene una única solución en $[0,1]$.
\end{ejercicio}


\begin{ejercicio}[Lectura 11, ejercicio 3]\label{ej:lectura-11-3}
Pruebe que $f(x)=x^2+x^{-1}$ no tiene puntos fijos en $[1,\infty)$.
\end{ejercicio}


\begin{ejercicio}[Lectura 11, ejercicio 4]\label{ej:lectura-11-4}
Sean $0<r<1$ y $\alpha\in\R$. En $C([0,r])$ defina
\[
 (\Phi f)(x)=\alpha+\int_0^x f(t)\,dt.
\]
Pruebe que $\Phi$ es contracción para la norma uniforme y encuentre su punto fijo.
\end{ejercicio}


\begin{ejercicio}[Lectura 11, ejercicio 5]\label{ej:lectura-11-5}\label{ch11:hiper-completo}
Demuestre que $(\mathcal F(X),H)$ es completo si y sólo si $X$ es completo.
\end{ejercicio}


\begin{ejercicio}[Lectura 11, ejercicio 6]\label{ej:lectura-11-6}
Calcule por iteración la solución del sistema
\[
 5x^2-y^2=0,\qquad y-\tfrac14(\sin x+\cos y)=0
\]
que se encuentra en la región positiva próxima a $(1/4,1/4)$.
\end{ejercicio}


\section{Existencia por densidad y por aproximaciones sucesivas}
Las bolas anidadas organizan dos formas de obtener existencia en espacios completos. En Baire, permiten conservar una cantidad numerable de condiciones densas; en el principio de contracción, las estimaciones geométricas conducen a un único punto fijo. Para iterar, la aplicación debe conservar el dominio, ese dominio debe ser completo y la constante de contracción debe ser estrictamente menor que uno. Verificados esos tres puntos, la misma demostración produce soluciones de ecuaciones, funciones definidas por ecuaciones integrales y conjuntos geométricos autosimilares, con cotas de error explícitas.
